\title[Israel 1: Abraham--Éxodo]{\bfseries{Israel, el pueblo de Dios}}
\subtitle{Estudio 1. De Abraham al Éxodo}
\author{}\institute{}\date{}
\titlegraphic{\includegraphics[width=0.85\linewidth,height=0.38\textheight,keepaspectratio]{portada.jpg}\\[-1pt]
  {\fontsize{5}{6}\selectfont\color{sepia}\textit{Nota.} Elaboración propia: Gn 15:5 en texto hebreo consonántico (BHS).}}
\AtBeginSection[]{\begin{frame}\frametitle{Contenidos}\tableofcontents[currentsection]\end{frame}}
\begin{document}
\frame{\titlepage}
\begin{frame}\frametitle{Contenidos}\tableofcontents\end{frame}

% ===================================================
\section{Introducción}
\begin{frame}{¿Por qué un pueblo?}
\begin{columns}[t,onlytextwidth]
\begin{column}{0.52\textwidth}
\begin{cita}Son israelitas; de ellos son la adopción, la gloria, los pactos, la promulgación de la ley, el culto y las promesas; de ellos son los padres, y de ellos, según la carne, vino el Mesías.\fuente{(Ro 9:4--5)}\end{cita}
\begin{cita}Porque irrevocables son los dones y el llamamiento de Dios.\fuente{(Ro 11:29)}\end{cita}
\end{column}
\begin{column}{0.45\textwidth}\small
Para entender por qué Pablo dice esto hay que volver al principio.\par\medskip
\destacar{Tesis 1.} Israel es creación de Dios: nace de una pareja estéril y de una promesa (Gn 11:30; Hch 7:5).\par\medskip
\destacar{Tesis 2.} La restauración está escrita desde el origen: antes de que exista el pueblo, Dios anuncia destierro y regreso (Gn 15:13--16).
\end{column}
\end{columns}
\end{frame}

\begin{frame}{Los ejes de la serie en esta etapa}
\centering
\begin{tikzpicture}
 \foreach \x/\n/\t/\d/\r in {0/1/{Nación santa}/{¿Para qué fue llamado Abraham y qué función tiene su descendencia ante las naciones?}/{Gn 12:1--3; 18:18--19; Éx 4:22; 19:4--6},
   4.2/2/{Restauración}/{¿Dónde aparece por primera vez el patrón de exilio y regreso, y en qué se funda?}/{Gn 15:13--16; 46:4; 50:24; Éx 2:24},
   8.4/3/{Proyección}/{¿Cómo releen Isaías, Jeremías y Pablo a Abraham y al Éxodo?}/{Is 51:1--3; Jer 31:35--37; Ro 11:28--29}}{
  \node[caja,anchor=north] (b\n) at (\x,0) {\parbox[t][3.2cm][t]{3.55cm}{\RaggedRight
    {\small\textbf{\textcolor{ladrillo}{\t}}}\\[5pt]\small \d\par\vfill{\scriptsize\color{sepia}(\r)}}};
  \node[circle,fill=cafe,text=piedra,font=\scriptsize\bfseries,inner sep=1.6pt] at (b\n.north west) {\n};}
\end{tikzpicture}\par\vspace{6pt}
{\scriptsize\color{sepia}Las citas bíblicas son traducción propia desde el texto hebreo (BHS) y griego (NA28). Numeración española [hebrea] cuando difiere.}
\end{frame}

% ===================================================
\section{Cronología}
\begin{frame}{Secuencia de la etapa}
\centering
\begin{tikzpicture}[x=0.78cm]
\draw[line width=1.6pt,cafe,-{Latex[length=3mm]}] (0,0) -- (15.7,0);
\foreach \x/\t/\b in {0.6/{Llamado de Abram}/{Gn 12:1--3},
                      2.7/{Pacto entre las piezas}/{Gn 15},
                      4.8/{Isaac, hijo de la promesa}/{Gn 21:1--7},
                      6.9/{Jacob es llamado Israel}/{Gn 32:28 [32:29]},
                      9.0/{Descenso a Egipto}/{Gn 46:1--7},
                      11.1/{Opresión y clamor}/{Éx 1--2},
                      13.2/{Pascua y salida}/{Éx 12--15},
                      15.1/{Sinaí}/{Éx 19 (Est. 2)}}{
  \fill[oro] (\x,0) circle (4pt);
  \node[above,font=\scriptsize\bfseries,text=ladrillo,text width=1.6cm,align=center] at (\x,0.2) {\t};
  \node[below,font=\tiny,text=sepia,text width=1.6cm,align=center] at (\x,-0.2) {\b};}
\draw[decorate,decoration={brace,amplitude=5pt,mirror},sepia] (0.4,-0.9) -- (7.1,-0.9) node[midway,below=5pt,font=\scriptsize\itshape]{Patriarcas en Canaán: la promesa};
\draw[decorate,decoration={brace,amplitude=5pt,mirror},sepia] (8.8,-0.9) -- (13.4,-0.9) node[midway,below=5pt,font=\scriptsize\itshape]{Egipto: el primer exilio y su fin};
\end{tikzpicture}\par\vspace{4pt}
{\fontsize{5}{6}\selectfont\color{sepia}\raggedleft\textit{Nota.} Elaboración propia. Secuencia narrativa; las fechas absolutas están en debate.\par}
\end{frame}

\begin{frame}{Fechas en debate}
\begin{columns}[t,onlytextwidth]
\begin{column}{0.48\textwidth}\small
\textbf{\textcolor{cafe}{¿Cuándo fue el Éxodo?}}
\begin{itemize}
\item \destacar{Datación temprana} (c. 1446 a.\,e.\,c.): 1~R 6:1 cuenta 480 años hasta el 4.º año de Salomón.
\item \destacar{Datación tardía} (s. \textsc{xiii} a.\,e.\,c.): la ciudad de Ramsés (Éx 1:11) se identifica con Pi-Ramsés, de Ramsés~II.
\item La \textbf{Estela de Merneptah} (c. 1208) ya menciona a Israel como pueblo en Canaán.
\end{itemize}
\end{column}
\begin{column}{0.50\textwidth}
\textbf{\small\textcolor{cafe}{Los 430 años de Éx 12:40}}\par\smallskip
\scriptsize\renewcommand{\arraystretch}{1.2}
\begin{tabular}{>{\RaggedRight\arraybackslash}p{2.0cm}>{\RaggedRight\arraybackslash}p{3.7cm}}
\toprule
\textbf{Fuente} & \textbf{Lectura}\\\midrule
Texto Masorético & 430 años en Egipto\\
LXX y Pent. samaritano & 430 años «en Egipto \emph{y en Canaán}»\\
Pablo (Gá 3:17) & 430 años de la promesa a la ley\\
Josefo (\emph{Ant.} 2.15.2) & 215 años en Egipto\\
Seder Olam; Rashi & 400 años desde Isaac; c. 210 en Egipto\\
\bottomrule
\end{tabular}
\end{column}
\end{columns}
\end{frame}

% ===================================================
\section{Arqueología}
\begin{frame}{El camino de Abraham}
\begin{columns}[c,onlytextwidth]
\begin{column}{0.70\textwidth}\centering\includegraphics[width=\linewidth,height=0.86\textheight,keepaspectratio]{mapa_patriarcas.pdf}\par\vspace{1pt}{\fontsize{4.5}{5.5}\selectfont\color{sepia}\textit{Nota.} Elaboración propia (GSHHS/WDB, Basemap); ubicaciones aproximadas.\par}\end{column}
\begin{column}{0.28\textwidth}\scriptsize\RaggedRight
{\scriptsize\RaggedRight
\tikz[baseline=-0.6ex]\draw[ladrillo,line width=1.5pt,line cap=round](0,0)--(0.6,0);~Taré y Abram: Ur $\to$ Harán (Gn 11:31)\par
\tikz[baseline=-0.6ex]\draw[ladrillo,line width=1.5pt,line cap=round](0,0)--(0.6,0);~Abram: Harán $\to$ Siquem (Gn 12:4–6)\par
\tikz[baseline=-0.6ex]\draw[oro,line width=1.5pt,line cap=round,dash pattern=on 3.4pt off 1.8pt](0,0)--(0.6,0);~Abram a Egipto (Gn 12:10)\par
\tikz[baseline=-0.6ex]\draw[cafe,line width=1.5pt,line cap=round,dash pattern=on 1.2pt off 1.4pt](0,0)--(0.6,0);~Jacob y su casa a Egipto (Gn 46)\par}
\medskip

\textbf{\textcolor{ladrillo}{Ur $\to$ Harán $\to$ Siquem}}: unos 1.500 km siguiendo el Éufrates y la ruta de Damasco (Gn 11:31--12:6).\par\medskip
Ur y Harán eran centros del culto a la luna: «vuestros padres [\ldots] servían a otros dioses» (Jos 24:2). El llamado es una ruptura religiosa, no solo geográfica.\par\medskip
Canaán queda entre los dos imperios del río: Egipto y Mesopotamia. Por allí pasarán Abram, Jacob y todo Israel.
\end{column}
\end{columns}
\end{frame}

\begin{frame}{Semitas en Egipto}
\centering
\begin{tikzpicture}
 \foreach \x/\n/\t/\d/\r in {0/1/{Beni Hasan}/{Pintura de la tumba de Jnumhotep~II: 37 «asiáticos» (\emph{aamu}) guiados por Abisha llegan a Egipto.}/{c. 1890 a.\,e.\,c.; Newberry, 1893},
   3.15/2/{Avaris}/{Tell el-Dab\textquoteright a, delta oriental (zona de Gosén): población cananea desde c. 1800 a.\,e.\,c.}/{Bietak, 1996},
   6.3/3/{Papiro Brooklyn}/{Lista de siervos: casi la mitad con nombres semíticos; una mujer llamada \emph{Shifra} (cf. Éx 1:15).}/{Hayes, 1955; Albright, 1954},
   9.45/4/{Merneptah}/{Estela: «Israel está devastado; su descendencia no existe». Primera mención de Israel fuera de la Biblia.}/{c. 1208 a.\,e.\,c.}}{
  \node[caja,anchor=north] (b\n) at (\x,0) {\parbox[t][3.35cm][t]{2.65cm}{\RaggedRight
    {\small\textbf{\textcolor{ladrillo}{\t}}}\\[4pt]\scriptsize \d\par\vfill{\tiny\color{sepia}(\r)}}};
  \node[circle,fill=cafe,text=piedra,font=\scriptsize\bfseries,inner sep=1.6pt] at (b\n.north west) {\n};}
\end{tikzpicture}\par\vspace{5pt}
\begin{minipage}{0.95\linewidth}\scriptsize Ninguno de estos hallazgos menciona a José o a Jacob: muestran que el mundo que describe Génesis es históricamente verosímil.\end{minipage}
\end{frame}

\begin{frame}{Lo que la arqueología no dice}
\begin{columns}[t,onlytextwidth]
\begin{column}{0.49\textwidth}\small
\begin{itemize}
\item \textbf{Paralelos de Nuzi y Mari.} Se usaron para fechar a los patriarcas; Thompson (1974) y Van Seters (1975) los criticaron; Kitchen (2003) defiende un marco del segundo milenio.
\item \textbf{Silencio egipcio.} Ninguna fuente egipcia narra el Éxodo. Las inscripciones reales no registraban derrotas, y el delta conserva mal los papiros (Hoffmeier, 1997).
\end{itemize}
\end{column}
\begin{column}{0.49\textwidth}\small
\begin{itemize}
\item \textbf{Camellos (Gn 12:16).} Restos del Aravá sitúan el camello doméstico en el sur del Levante hacia el s. \textsc{x} a.\,e.\,c. (Sapir-Hen y Ben-Yosef, 2013); otros apuntan a un uso anterior y limitado.
\end{itemize}\vspace{6pt}
\begin{tcolorbox}[colback=piedra,colframe=cafe,boxrule=0.8pt,arc=3pt,fontupper=\small]
La ausencia de evidencia directa obliga a la cautela; no autoriza a negar el relato.
\end{tcolorbox}
\end{column}
\end{columns}
\end{frame}

% ===================================================
\section{El relato}
\begin{frame}{Un nombre recibido}
\fotolado{panel_gn12.jpg}{Elaboración propia: Gn 12:2--3 en texto hebreo consonántico (BHS).}
\begin{columns}[c,onlytextwidth]
\begin{column}{0.57\textwidth}
\small Babel: «Hagámonos un nombre» (\heb{וְנַעֲשֶׂה־לָּנוּ שֵׁם}, Gn 11:4). Dios responde con un hombre:
\begin{cita}Haré de ti una nación grande, te bendeciré y engrandeceré tu nombre, y sé tú una bendición [\ldots]; y en ti serán benditas todas las familias de la tierra.\fuente{(Gn 12:2--3)}\end{cita}
\scriptsize
\begin{itemize}
\item La iniciativa es de Dios: «haré», «bendeciré», «engrandeceré».
\item La nación (\heb{גּוֹי}, \emph{goy}) recibe una tarea: \heb{וֶהְיֵה בְּרָכָה} (\emph{ve-heyé berajá}).
\item La meta es universal: las familias de Gn 10, dispersadas en Gn 11.
\end{itemize}
\end{column}
\begin{column}{0.41\textwidth}\end{column}
\end{columns}
\end{frame}

\begin{frame}{«En ti serán benditas»: del hebreo a Pablo}
\centering
\begin{tikzpicture}[node distance=0.9cm]
\node[caja,text width=3.3cm,align=center] (h) {\textbf{Hebreo (Gn 12:3)}\\[3pt]\heb{וְנִבְרְכוּ בְךָ}\\[2pt]\emph{ve-nivrejú vejá}\\[2pt]nifal: pasivo o reflexivo};
\node[caja,text width=3.3cm,align=center,right=of h] (g) {\textbf{LXX}\\[3pt]\grc{ἐνευλογηθήσονται ἐν σοί}\\[2pt]futuro pasivo:\\«serán benditas»};
\node[caja,text width=3.3cm,align=center,right=of g] (p) {\textbf{Pablo (Gá 3:8)}\\[3pt]\grc{πάντα τὰ ἔθνη}\\[2pt]«todas las naciones»\\(de Gn 18:18; 22:18)};
\draw[-{Latex[length=2.5mm]},oro!90!ladrillo,line width=1.4pt] (h) -- (g);
\draw[-{Latex[length=2.5mm]},oro!90!ladrillo,line width=1.4pt] (g) -- (p);
\end{tikzpicture}\par\vspace{8pt}
{\small\itshape\color{sepia}Pablo funde Gn 12:3 con Gn 18:18 y llama a esta promesa «evangelio anunciado de antemano a Abraham».}\par\vspace{6pt}
\begin{minipage}{0.92\linewidth}\precision{Gn 22:18 y 26:4 usan el hitpael \heb{וְהִתְבָּרְכוּ}, que favorece «se bendecirán». Ambas lecturas ligan el destino de las naciones a Abraham.}\end{minipage}
\end{frame}

\begin{frame}{El pacto entre las piezas (Gn 15)}
\begin{columns}[t,onlytextwidth]
\begin{column}{0.5\textwidth}
\centering{\large\heb{וְהֶאֱמִן בַּיהוָה וַיַּחְשְׁבֶהָ לּוֹ צְדָקָה}}\par\smallskip
\begin{cita}Y creyó a YHWH, y él se lo contó por justicia.\fuente{(Gn 15:6)}\end{cita}
\small \heb{הֶאֱמִן} (\emph{he'emín}), raíz \emph{'mn} (\emph{amén}, \emph{emuná}): apoyarse en algo firme. Base de Ro 4 y Gá 3.
\end{column}
\begin{column}{0.48\textwidth}\small
\begin{itemize}
\item Abram parte los animales; \textbf{solo Dios pasa} entre las mitades como «horno humeante y antorcha» (Gn 15:17).
\item En los tratados antiguos, quien pasaba se maldecía a sí mismo si rompía el pacto (Jer 34:18--20).
\item Dios asume solo la obligación. Tradición judía: \emph{berit bein ha-betarim}.
\end{itemize}
\end{column}
\end{columns}
\end{frame}

\begin{frame}{El primer anuncio de exilio y regreso}
\begin{cita}Tu descendencia será extranjera en una tierra que no es suya; los harán servir y los oprimirán cuatrocientos años. Pero a la nación a la que servirán la juzgaré yo, y después saldrán con gran riqueza [\ldots]. Y en la cuarta generación \textbf{volverán} acá.\fuente{(Gn 15:13--14, 16)}\end{cita}
\centering
\begin{tikzpicture}[node distance=0.55cm]
\foreach \i/\t/\s [count=\k] in {0/Promesa/{Gn 15:5--7},1/Destierro/{15:13},2/Juicio al opresor/{15:14a},3/Regreso/{15:14b, 16}}{
  \node[caja,text width=2.3cm,align=center,minimum height=1.0cm] (n\k) at (\i*3.1,0) {\textbf{\textcolor{ladrillo}{\t}}\\[1pt]{\tiny\color{sepia}\s}};}
\foreach \a/\b in {1/2,2/3,3/4}{\draw[-{Latex[length=2.2mm]},oro!90!ladrillo,line width=1.3pt] (n\a) -- (n\b);}
\end{tikzpicture}\par\vspace{4pt}
\small \heb{יָשׁוּבוּ} (\emph{yashuvu}), de \heb{שׁוּב} (\emph{shuv}): verbo del regreso y del arrepentimiento (Dt 30:3; Jer 29:14).
\end{frame}

\begin{frame}{Pacto perpetuo y Akedá}
\begin{columns}[t,onlytextwidth]
\begin{column}{0.49\textwidth}
\textbf{\textcolor{cafe}{Gn 17: el pacto perpetuo}}
\begin{cita}Estableceré mi pacto [\ldots] con tu descendencia después de ti, por sus generaciones, como pacto perpetuo (\upshape\heb{לִבְרִית עוֹלָם}).\fuente{(Gn 17:7)}\end{cita}
\small Abram pasa a ser Abraham: \heb{אַב־הֲמוֹן גּוֹיִם}, «padre de multitud de naciones» (Gn 17:5; Ro 4:17).
\end{column}
\begin{column}{0.49\textwidth}
\textbf{\textcolor{cafe}{Gn 22: la promesa entregada}}\par\smallskip
\small
\begin{itemize}
\item Isaac, hijo de la imposibilidad (Gn 18:11--14), es ofrecido.
\item \heb{יְהוָה יִרְאֶה} (\emph{YHWH yiré}): «YHWH proveerá» (22:14).
\item Promesa jurada: «en tu descendencia serán benditas todas las naciones» (22:18).
\item Moriah y el monte del templo (2 Cr 3:1); figura de resurrección (Heb 11:17--19).
\end{itemize}
\end{column}
\end{columns}
\end{frame}

\begin{frame}{Jacob llamado Israel}
\begin{columns}[t,onlytextwidth]
\begin{column}{0.5\textwidth}
\begin{cita}No se dirá más tu nombre Jacob, sino Israel, porque has luchado con Dios y con los hombres, y has prevalecido.\fuente{(Gn 32:28 [32:29])}\end{cita}
\small El texto explica el nombre con \emph{sará}, «luchar». Es etimología por sonido; el sentido gramatical probable es «Dios lucha» o «que Dios prevalezca» (Sarna, 1989).
\end{column}
\begin{column}{0.48\textwidth}\small
\textbf{\textcolor{cafe}{Dos expresiones para guardar}}
\begin{itemize}
\item \heb{גּוֹי וּקְהַל גּוֹיִם} (\emph{goy u-qehal goyim}): «una nación y un conjunto de naciones» (Gn 35:11).
\item \heb{מְלֹא־הַגּוֹיִם} (\emph{melo ha-goyim}): Efraín será «plenitud de las naciones» (Gn 48:19).
\end{itemize}
\precision{volveremos a ellas en la división del reino y en Romanos 11 («plenitud de los gentiles», Ro 11:25).}
\end{column}
\end{columns}
\end{frame}

\begin{frame}{José y el primer remanente}
\begin{cita}Dios me envió delante de vosotros para preservaros un \textbf{remanente} en la tierra y para daros vida por medio de una gran liberación.\fuente{(Gn 45:7)}\end{cita}
\centering
\begin{tikzpicture}[node distance=0.8cm]
\node[caja,text width=3.2cm,align=center] (a) {\textbf{Génesis}\\[2pt]\heb{שְׁאֵרִית} (\emph{she'erit})\\{\tiny Gn 45:7: la familia sobrevive}};
\node[caja,text width=3.2cm,align=center,right=of a] (b) {\textbf{Isaías}\\[2pt]\heb{שְׁאָר יָשׁוּב}\\{\tiny Is 10:20--22: «un remanente volverá»}};
\node[caja,text width=3.2cm,align=center,right=of b] (c) {\textbf{Pablo}\\[2pt]\grc{λεῖμμα κατ᾽ ἐκλογὴν χάριτος}\\{\tiny Ro 11:5: «remanente escogido por gracia»}};
\draw[-{Latex[length=2.5mm]},oro!90!ladrillo,line width=1.4pt] (a) -- (b);
\draw[-{Latex[length=2.5mm]},oro!90!ladrillo,line width=1.4pt] (b) -- (c);
\end{tikzpicture}\par\vspace{6pt}
{\small\itshape\color{sepia}La primera vez que Israel es remanente, lo es para sobrevivir y para que otros vivan: José alimenta a «toda la tierra» (Gn 41:57).}
\end{frame}

\begin{frame}{Dios baja al exilio y promete visitar}
\begin{columns}[t,onlytextwidth]
\begin{column}{0.49\textwidth}
\centering{\large\heb{אָנֹכִי אֵרֵד עִמְּךָ מִצְרַיְמָה וְאָנֹכִי אַעַלְךָ גַם־עָלֹה}}\par\smallskip
\begin{cita}Yo descenderé contigo a Egipto, y yo también te haré subir de nuevo.\fuente{(Gn 46:4)}\end{cita}
\small El exilio no separa a Israel de su Dios.
\end{column}
\begin{column}{0.49\textwidth}\small
\textbf{\textcolor{cafe}{La fórmula de la visitación}}
\begin{itemize}
\item José: \heb{פָּקֹד יִפְקֹד} (\emph{paqod yifqod}), «Dios ciertamente os visitará» (Gn 50:24--25).
\item Moisés en la zarza: \heb{פָּקֹד פָּקַדְתִּי} (Éx 3:16).
\item Los huesos de José suben con Israel (Éx 13:19; Heb 11:22).
\item Zacarías: Dios «ha visitado y redimido a su pueblo» (Lc 1:68).
\end{itemize}
\precision{una tradición midrásica dice que la fórmula era la señal de redención transmitida desde José (Pirqé de-Rabí Eliézer 48).}
\end{column}
\end{columns}
\end{frame}

\begin{frame}{El Éxodo: Dios se acuerda}
\begin{columns}[t,onlytextwidth]
\begin{column}{0.47\textwidth}
\begin{cita}Y oyó Dios su gemido, y se acordó Dios de su pacto con Abraham, con Isaac y con Jacob.\fuente{(Éx 2:24)}\end{cita}
\small \heb{זָכַר} (\emph{zajar}): no recuperar un dato olvidado, sino actuar en virtud de un compromiso.\par\smallskip
\heb{בְּנִי בְכֹרִי יִשְׂרָאֵל}: «Israel es mi hijo, mi primogénito» (Éx 4:22).\par\smallskip
\heb{עֵרֶב רַב}: con Israel sale una «multitud mixta» (Éx 12:38, 48--49).
\end{column}
\begin{column}{0.51\textwidth}
\centering\textbf{\small\textcolor{cafe}{Los verbos de la redención (Éx 6:6--8)}}\par\smallskip
\begin{tikzpicture}[y=-0.62cm]
\foreach \i/\v/\h in {0/Os sacaré/{וְהוֹצֵאתִי},1/Os libraré/{וְהִצַּלְתִּי},2/Os redimiré/{וְגָאַלְתִּי},3/Os tomaré por pueblo/{וְלָקַחְתִּי},4/Os introduciré en la tierra/{וְהֵבֵאתִי}}{
  \node[caja,text width=4.6cm,minimum height=0.48cm,inner sep=2.5pt] at (0,\i) {\textbf{\v}\hfill\heb{\h}};}
\draw[decorate,decoration={brace,amplitude=4pt},sepia] (2.6,-0.25) -- (2.6,2.25) node[midway,right=4pt,font=\tiny,text width=1.3cm]{Pascua, mar, Sinaí};
\draw[decorate,decoration={brace,amplitude=4pt},sepia] (2.6,3.75) -- (2.6,4.25) node[midway,right=4pt,font=\tiny,text width=1.3cm]{Estudio 2};
\end{tikzpicture}
\end{column}
\end{columns}
\end{frame}

% ===================================================
\section{Vocación}
\begin{frame}{Un pueblo sacerdotal}
\fotolado{panel_ex.jpg}{Elaboración propia: Éx 4:22 y 19:6 en texto hebreo consonántico (BHS).}
\begin{columns}[c,onlytextwidth]
\begin{column}{0.57\textwidth}
\begin{cita}Vosotros me seréis un reino de sacerdotes y una nación santa.\fuente{(Éx 19:6)}\end{cita}
\small \heb{מַמְלֶכֶת כֹּהֲנִים וְגוֹי קָדוֹשׁ}. La fórmula se dice en el Sinaí (Estudio 2), pero su contenido ya está en los patriarcas:
\scriptsize
\begin{itemize}
\item \textbf{Altar}: Abraham invoca el nombre (Gn 12:7--8; 13:18).
\item \textbf{Intercesión}: por Sodoma y por Abimelec (Gn 18:22--33; 20:7).
\item \textbf{Enseñanza}: «justicia y juicio» a sus hijos (Gn 18:19).
\item \textbf{Preservar la vida}: José y la hambruna (Gn 41:57).
\item \textbf{Primogénito}: representa a los hermanos (Éx 4:22; Nm 3:12).
\end{itemize}
\end{column}
\begin{column}{0.41\textwidth}\end{column}
\end{columns}
\end{frame}

\begin{frame}{El llamamiento de Israel}
\centering
\begin{tikzpicture}
 \foreach \x/\y/\n/\t/\r in {0/0/1/{Reino de sacerdotes y nación santa}/{Éx 19:6},
    4.2/0/2/{Luz para las naciones}/{Is 49:6},
    8.4/0/3/{Custodio de la palabra de Dios}/{Dt 4:5--8; Sal 147:19--20; Ro 3:1--2},
    2.1/-2.35/4/{Testigo de la fidelidad de Dios}/{Is 43:10--12},
    6.3/-2.35/5/{Heredero de la tierra}/{Gn 13:14--15; 15:18}}{
  \node[caja,anchor=north] (b\n) at (\x,\y) {\parbox[t][1.55cm][t]{3.55cm}{\RaggedRight
    {\small\textbf{\textcolor{ladrillo}{\t}}}\par\vfill{\scriptsize\color{sepia}(\r)}}};
  \node[circle,fill=cafe,text=piedra,font=\scriptsize\bfseries,inner sep=1.6pt] at (b\n.north west) {\n};}
\end{tikzpicture}\par\vspace{6pt}
\begin{minipage}{0.95\linewidth}\scriptsize Esta etapa contiene en germen los cinco. Isaías los proyecta a la restauración: «seréis llamados sacerdotes de YHWH» (Is 61:6), y de las naciones tomará «sacerdotes y levitas» (Is 66:21).\par
{\fontsize{5}{6}\selectfont\color{sepia}\textit{Nota.} Adaptado de Departamento de Educación (s.\,f.), \emph{Restauración de Israel}, Clase 7.}\end{minipage}
\end{frame}

% ===================================================
\section{Tradición judía}
\begin{frame}{Abraham en el midrash}
\small Narraciones de la tradición, no datos históricos: muestran cómo Israel leyó su origen.\par\medskip
\centering
\begin{tikzpicture}
 \foreach \x/\n/\t/\d/\r in {0/1/{El rompedor de ídolos}/{Abram rompe las imágenes del taller de su padre Taré. Da forma narrativa a Jos 24:2.}/{Bereshit Rabbah 38:13},
   4.2/2/{«Abram el hebreo»}/{\heb{הָעִבְרִי} (\emph{ha-ivrí}) leído desde \emph{éver}, «el otro lado»: el mundo de un lado y Abram del otro.}/{Bereshit Rabbah 42:8},
   8.4/3/{Las diez pruebas}/{«Con diez pruebas fue probado Abraham, nuestro padre, y las superó todas».}/{Mishná, Avot 5:3}}{
  \node[caja,anchor=north] (b\n) at (\x,0) {\parbox[t][2.75cm][t]{3.55cm}{\RaggedRight
    {\small\textbf{\textcolor{ladrillo}{\t}}}\\[5pt]\small \d\par\vfill{\scriptsize\color{sepia}(\r)}}};
  \node[circle,fill=cafe,text=piedra,font=\scriptsize\bfseries,inner sep=1.6pt] at (b\n.north west) {\n};}
\end{tikzpicture}
\end{frame}

\begin{frame}{Los padres en la oración diaria}
\centering{\large\heb{וְזוֹכֵר חַסְדֵי אָבוֹת וּמֵבִיא גּוֹאֵל לִבְנֵי בְנֵיהֶם}}\par\smallskip
\begin{cita}Dios de Abraham, Dios de Isaac y Dios de Jacob [\ldots] que se acuerda de las bondades de los padres y trae un redentor a los hijos de sus hijos.\fuente{(Amidá, bendición \emph{Avot})}\end{cita}
\begin{columns}[t,onlytextwidth]
\begin{column}{0.49\textwidth}\small
\textbf{\textcolor{cafe}{Tradición judía}}: \emph{zejut avot}, «el mérito de los padres». Se ora tres veces al día.
\end{column}
\begin{column}{0.49\textwidth}\small
\textbf{\textcolor{cafe}{Pablo}}: Israel es «amado por causa de los padres» (Ro 11:28). El fundamento no es el mérito de los padres, sino la fidelidad de Dios a lo que les prometió.
\end{column}
\end{columns}
\end{frame}

\begin{frame}{Memoria que se vuelve esperanza}
\begin{columns}[t,onlytextwidth]
\begin{column}{0.32\textwidth}\small
\textbf{\textcolor{ladrillo}{Rosh Hashaná}}\par\smallskip
Se leen Gn 21--22 y se toca un cuerno de carnero: «para que me acuerde de vosotros de la atadura de Isaac» (Talmud, Rosh Hashaná 16a).
\end{column}
\begin{column}{0.33\textwidth}\small
\textbf{\textcolor{ladrillo}{Pésaj}}\par\smallskip
«En cada generación, cada uno debe verse como si hubiera salido de Egipto» (Mishná, Pesajim 10:5). Cuatro copas por los verbos de Éx 6:6--7 (T. de Jerusalén, Pesajim 10:1). Final: «el año próximo en Jerusalén».
\end{column}
\begin{column}{0.32\textwidth}\small
\textbf{\textcolor{ladrillo}{La Shejiná en el exilio}}\par\smallskip
«A todo lugar adonde fueron desterrados, la Shejiná fue con ellos: a Egipto [\ldots], a Babilonia» (Talmud, Meguilá 29a). Desarrolla Gn 46:4.
\end{column}
\end{columns}
\end{frame}

% ===================================================
\section{Restauración}
\begin{frame}{El patrón en la Torá}
\centering\scriptsize\renewcommand{\arraystretch}{1.25}
\begin{tabular}{>{\RaggedRight\arraybackslash}p{1.9cm}>{\RaggedRight\arraybackslash}p{4.6cm}>{\RaggedRight\arraybackslash}p{4.4cm}}
\toprule
\textbf{Texto} & \textbf{Promesa} & \textbf{Proyección}\\\midrule
Gn 15:13--16 & Destierro, juicio al opresor, regreso & Dt 30:1--5; Jer 29:10--14\\
Gn 17:7 & Pacto perpetuo con la descendencia & Jer 31:35--37 [31:34--36]; Ro 11:29\\
Gn 45:7 & Un remanente preservado & Is 10:20--22; Ro 11:5\\
Gn 46:4 & Dios baja al exilio y hace subir & Is 43:2; Ez 11:16\\
Gn 50:24 & «Dios ciertamente os visitará» & Éx 3:16; Lc 1:68\\
Éx 2:24 & Dios se acuerda del pacto & Lv 26:42--45; Lc 1:72--73\\
Éx 6:6--8 & Sacar, librar, redimir, tomar, introducir & Is 43:1; Jer 16:14--15\\
\bottomrule
\end{tabular}\par\vspace{6pt}
\begin{cita}Me acordaré de mi pacto con Jacob, y también de mi pacto con Isaac, y también de mi pacto con Abraham me acordaré, y de la tierra me acordaré.\fuente{(Lv 26:42)}\end{cita}
\end{frame}

\begin{frame}{Isaías: mirad la roca}
\fotolado{panel_is51.jpg}{Elaboración propia: Is 51:2 en texto hebreo consonántico (BHS).}
\begin{columns}[c,onlytextwidth]
\begin{column}{0.57\textwidth}
\begin{cita}Mirad a Abraham vuestro padre y a Sara que os dio a luz; porque era uno solo cuando lo llamé, y lo bendije y lo multipliqué. Ciertamente consolará YHWH a Sion.\fuente{(Is 51:2--3)}\end{cita}
\scriptsize
\begin{itemize}
\item \textbf{El origen prueba la restauración}: si Dios hizo una nación de uno solo, puede rehacer a Sion.
\item \textbf{Is 41:8--9}: «descendencia de Abraham mi amigo [\ldots] te escogí y no te deseché».
\item \textbf{Is 43:19}: \heb{הִנְנִי עֹשֶׂה חֲדָשָׁה}, «hago una cosa nueva»: el nuevo éxodo.
\item \textbf{Is 11:11, 16}: el Señor extenderá su mano «por segunda vez» (\heb{שֵׁנִית}), como cuando Israel subió de Egipto.
\end{itemize}
\end{column}
\begin{column}{0.41\textwidth}\end{column}
\end{columns}
\end{frame}

\begin{frame}{Jeremías: un pacto tan firme como el cosmos}
\begin{columns}[t,onlytextwidth]
\begin{column}{0.49\textwidth}\small
\begin{itemize}
\item \textbf{Jer 16:14--15; 23:7--8}: ya no se jurará por el Dios que sacó a Israel de Egipto, sino por el que lo trajo «de todas las tierras adonde los había arrojado».
\item \textbf{Jer 31:9 [31:8]}: «soy a Israel por padre, y Efraín es mi primogénito» (cf. Éx 4:22).
\end{itemize}
\end{column}
\begin{column}{0.49\textwidth}
\begin{cita}Si faltaren estas leyes delante de mí, dice YHWH, también la descendencia de Israel faltará para no ser nación delante de mí para siempre.\fuente{(Jer 31:36 [31:35])}\end{cita}
\begin{cita}[fontupper=\scriptsize\itshape]Si no permanece mi pacto con el día y la noche [\ldots], también desecharé la descendencia de Jacob [\ldots]. Porque haré volver a sus cautivos y tendré de ellos misericordia.\fuente{(Jer 33:25--26)}\end{cita}
\end{column}
\end{columns}
\end{frame}

\begin{frame}{Miqueas y Lucas: el juramento a los padres}
\centering{\large\heb{תִּתֵּן אֱמֶת לְיַעֲקֹב חֶסֶד לְאַבְרָהָם}}\par\smallskip
\begin{cita}Darás fidelidad a Jacob, y misericordia a Abraham, como lo juraste a nuestros padres desde los días antiguos.\fuente{(Mi 7:20)}\end{cita}
\begin{columns}[t,onlytextwidth]
\begin{column}{0.49\textwidth}
\begin{cita}[fontupper=\scriptsize\itshape]Socorrió a Israel su siervo, acordándose de la misericordia, como habló a nuestros padres, a Abraham y a su descendencia para siempre.\fuente{(Lc 1:54--55)}\end{cita}
\end{column}
\begin{column}{0.49\textwidth}
\begin{cita}[fontupper=\scriptsize\itshape]Ha visitado y redimido a su pueblo [\ldots] para acordarse de su santo pacto, del juramento que hizo a Abraham nuestro padre.\fuente{(Lc 1:68, 72--73)}\end{cita}
\end{column}
\end{columns}
{\small\itshape\color{sepia}\centering Los verbos son los del Éxodo: visitar, redimir, acordarse.\par}
\end{frame}

\begin{frame}{Pablo: Ro 11:28--29}
\fotolado{panel_gal38.jpg}{Elaboración propia: Gá 3:8 en escritura uncial (texto NA28).}
\begin{columns}[c,onlytextwidth]
\begin{column}{0.57\textwidth}
\begin{cita}En cuanto a la elección, son amados por causa de los padres; porque irrevocables son los dones y el llamamiento de Dios.\fuente{(Ro 11:28--29)}\end{cita}
\scriptsize
\begin{itemize}
\item \grc{ἀμεταμέλητα}: «sin arrepentimiento»; Dios no se retracta.
\item \textbf{Gá 3:8}: la promesa a Abraham es «evangelio anunciado de antemano».
\item \textbf{Ro 4:17}: Abraham, padre de los creyentes de las naciones.
\item \textbf{Ro 11:16}: «si la raíz es santa, también las ramas» (raíz = los patriarcas, según muchos; Moo, 1996).
\item \textbf{Ro 15:8}: Cristo vino «para confirmar las promesas hechas a los padres».
\end{itemize}
\end{column}
\begin{column}{0.41\textwidth}\end{column}
\end{columns}
\end{frame}

% ===================================================
\section{Síntesis}
\begin{frame}{Síntesis}
\begin{tcolorbox}[colback=piedra,colframe=ladrillo,boxrule=1pt,arc=3pt,fontupper=\small]
\textbf{\textcolor{cafe}{Israel en su origen:}} un pueblo que Dios crea de una pareja estéril, llama a ser bendición para las naciones y sostiene con un pacto perpetuo que anuncia de antemano su destierro y su regreso. Por ese pacto Dios «se acordó» en Egipto (Éx 2:24), por él Isaías y Jeremías anuncian un nuevo éxodo, y por él Pablo afirma que Israel sigue siendo «amado por causa de los padres» (Ro 11:28).
\end{tcolorbox}
\scriptsize\renewcommand{\arraystretch}{1.2}
\begin{tabular}{>{\RaggedRight\arraybackslash}p{3.0cm}>{\RaggedRight\arraybackslash}p{8.3cm}}
\toprule
\textbf{Eje} & \textbf{Lo que deja esta etapa}\\\midrule
Creación de Dios & Sara estéril, Isaac hijo de la promesa (Gn 11:30; 21:1--7; Is 51:1--3)\\
Vocación sacerdotal & Bendición para las naciones; altar, intercesión, enseñanza, primogenitura\\
Restauración & Gn 15:13--16; 46:4; 50:24; Éx 2:24; Lv 26:42\\
Proyección & Isaías (nuevo éxodo), Jeremías (pacto como el cosmos), Pablo (Ro 11:28--29)\\
\bottomrule
\end{tabular}
\end{frame}

\begin{frame}{Hacia el Estudio 2}
\begin{columns}[c,onlytextwidth]
\begin{column}{0.55\textwidth}
\begin{cita}Vosotros visteis lo que hice a los egipcios, y cómo os tomé sobre alas de águila y os he traído a mí.\fuente{(Éx 19:4)}\end{cita}
\small El Éxodo no termina en el mar. En el Sinaí, la vocación vista en germen recibirá forma de \destacar{pacto}, \destacar{ley} y \destacar{santuario}.
\end{column}
\begin{column}{0.42\textwidth}
\begin{tcolorbox}[colback=cafe,colframe=cafeosc,coltext=piedra,boxrule=0.8pt,arc=3pt,fontupper=\small]
\textbf{Estudio 2}\par\smallskip
Del Éxodo a la muerte de Josué: Sinaí, tabernáculo, desierto, conquista y reparto de la tierra.
\end{tcolorbox}
\end{column}
\end{columns}
\end{frame}

% ===================================================
\section{Reflexión}
\begin{frame}{Para la reflexión}
\small
\begin{enumerate}
\item Isaías invita a «mirar la roca de donde fuisteis tallados» (Is 51:1). ¿Qué cambia al recordar que Israel nació de una mujer estéril?
\item Gn 15 anuncia el sufrimiento de Israel antes de que exista. ¿Qué dice esto sobre la relación entre los juicios de Dios y su fidelidad?
\item Si Abraham fue llamado a ser bendición para las naciones, ¿qué responsabilidad tienen hoy los creyentes de las naciones hacia Israel (Ro 11:18; 15:27)?
\item ¿En qué se parecen y en qué se diferencian el \emph{zejut avot} rabínico y la afirmación de Pablo en Ro 11:28?
\end{enumerate}
\end{frame}

% ===================================================
\section{Glosario}
\begin{frame}{Glosario}
\begin{columns}[t,onlytextwidth]
\begin{column}{0.49\textwidth}\scriptsize
\textbf{Akedá} (\heb{עֲקֵדָה}): «atadura»; relato de Gn 22.\par\smallskip
\textbf{Berit} (\heb{בְּרִית}): pacto; \emph{berit olam}, pacto perpetuo.\par\smallskip
\textbf{Berit bein ha-betarim}: «pacto entre las piezas» (Gn 15).\par\smallskip
\textbf{Bejor} (\heb{בְּכוֹר}): primogénito (Éx 4:22; Jer 31:9).\par\smallskip
\textbf{Emuná} (\heb{אֱמוּנָה}): firmeza, fidelidad, fe.\par\smallskip
\textbf{Érev rav} (\heb{עֵרֶב רַב}): multitud mixta (Éx 12:38).\par\smallskip
\textbf{Ha-ivrí} (\heb{הָעִבְרִי}): «el hebreo» (Gn 14:13).
\end{column}
\begin{column}{0.49\textwidth}\scriptsize
\textbf{Mamléjet kohaním} (\heb{מַמְלֶכֶת כֹּהֲנִים}): reino de sacerdotes (Éx 19:6).\par\smallskip
\textbf{Paqad} (\heb{פָּקַד}): visitar, intervenir a favor (Gn 50:24).\par\smallskip
\textbf{She'erit} (\heb{שְׁאֵרִית}): remanente (Gn 45:7).\par\smallskip
\textbf{Shejiná} (\heb{שְׁכִינָה}): presencia de Dios con su pueblo (término rabínico).\par\smallskip
\textbf{Shuv} (\heb{שׁוּב}): volver; regreso y arrepentimiento.\par\smallskip
\textbf{Zejut avot} (\heb{זְכוּת אָבוֹת}): mérito de los padres.
\end{column}
\end{columns}
\end{frame}

% ===================================================
\section{Referencias}
\begin{frame}[allowframebreaks]{Referencias}
\tiny
\refapa{Albright, W.~F. (1954). Northwest-Semitic names in a list of Egyptian slaves from the eighteenth century B.C. \emph{Journal of the American Oriental Society, 74}(4), 222--233.}
\refapa{\emph{Biblia Hebraica Stuttgartensia} (5.ª ed.). (1997). Deutsche Bibelgesellschaft.}
\refapa{Bietak, M. (1996). \emph{Avaris: The capital of the Hyksos}. British Museum Press.}
\refapa{Danby, H. (Trad.). (1933). \emph{The Mishnah}. Oxford University Press.}
\refapa{Departamento de Educación, Congregación Iglesia Israelita del Nuevo Pacto. (s.\,f.). \emph{Restauración de Israel} (Clase 7) [Presentación].}
\refapa{Epstein, I. (Ed.). (1935--1948). \emph{The Babylonian Talmud}. Soncino Press.}
\refapa{Freedman, H., \& Simon, M. (Eds.). (1939). \emph{Midrash Rabbah}. Soncino Press.}
\refapa{Friedlander, G. (Trad.). (1916). \emph{Pirkê de Rabbi Eliezer}. Kegan Paul.}
\refapa{Guggenheimer, H.~W. (1998). \emph{Seder Olam: The rabbinic view of biblical chronology}. Jason Aronson.}
\refapa{Hayes, W.~C. (1955). \emph{A papyrus of the late Middle Kingdom in the Brooklyn Museum}. Brooklyn Museum.}
\refapa{Hays, R.~B. (1989). \emph{Echoes of Scripture in the letters of Paul}. Yale University Press.}
\refapa{Hoffmeier, J.~K. (1997). \emph{Israel in Egypt: The evidence for the authenticity of the Exodus tradition}. Oxford University Press.}
\refapa{Josefo, F. (1987). \emph{The works of Josephus} (W. Whiston, Trad.). Hendrickson.}
\refapa{Kitchen, K.~A. (2003). \emph{On the reliability of the Old Testament}. Eerdmans.}
\refapa{Moo, D.~J. (1996). \emph{The epistle to the Romans}. Eerdmans.}
\refapa{Neusner, J. (Trad.). (1982--1994). \emph{The Talmud of the Land of Israel}. University of Chicago Press.}
\refapa{Newberry, P.~E. (1893). \emph{Beni Hasan} (Parte I). Egypt Exploration Fund.}
\refapa{\emph{Novum Testamentum Graece} (28.ª ed.). (2012). Deutsche Bibelgesellschaft.}
\refapa{ONE FOR ISRAEL Ministry. (2025, 3 de julio). \emph{What does «all Israel will be saved» mean in Romans 11?} \url{https://www.oneforisrael.org/what-does-all-israel-will-be-saved-mean-in-romans-11/}}
\refapa{Rahlfs, A., \& Hanhart, R. (Eds.). (2006). \emph{Septuaginta} (ed. rev.). Deutsche Bibelgesellschaft.}
\refapa{Sacks, J. (Trad.). (2009). \emph{The Koren Siddur}. Koren Publishers.}
\refapa{Sapir-Hen, L., \& Ben-Yosef, E. (2013). The introduction of domestic camels to the southern Levant: Evidence from the Aravah Valley. \emph{Tel Aviv, 40}(2), 277--285.}
\refapa{Sarna, N.~M. (1989). \emph{Genesis} (The JPS Torah Commentary). Jewish Publication Society.}
\refapa{Sarna, N.~M. (1991). \emph{Exodus} (The JPS Torah Commentary). Jewish Publication Society.}
\refapa{Thompson, T.~L. (1974). \emph{The historicity of the patriarchal narratives}. De Gruyter.}
\refapa{Van Seters, J. (1975). \emph{Abraham in history and tradition}. Yale University Press.}
\refapa{Wright, N.~T. (2013). \emph{Paul and the faithfulness of God}. Fortress Press.}
\end{frame}
\end{document}
